\documentclass[12pt, a4paper]{article}

% ── PACKAGES ──────────────────────────────────────────────────
\usepackage[utf8]{inputenc}
\usepackage[T1]{fontenc}
\usepackage{geometry}
\geometry{margin=2.5cm}
\usepackage{hyperref}
\hypersetup{
    colorlinks=true,
    linkcolor=blue,
    urlcolor=blue,
    citecolor=blue,
    pdftitle={The EZH2 Paradox in Triple-Negative Breast
              Cancer: Both Arms Confirmed in GSE25066 and
              Resolution via Waddington Attractor Geometry},
    pdfauthor={Eric Robert Lawson},
    pdfkeywords={EZH2, paradox, TNBC, triple-negative
                 breast cancer, pCR, pathological complete
                 response, distant recurrence free survival,
                 DRFS, neoadjuvant chemotherapy, GSE25066,
                 attractor geometry, Waddington landscape,
                 tazemetostat, EZH2 inhibitor, maintenance
                 therapy, attractor reconstitution,
                 OrganismCore}
}
\usepackage{amsmath}
\usepackage{booktabs}
\usepackage{longtable}
\usepackage{array}
\usepackage{parskip}
\usepackage{microtype}
\usepackage{authblk}
\usepackage{titlesec}
\usepackage{fancyhdr}
\usepackage{xcolor}
\usepackage{float}
\usepackage{enumitem}
\usepackage{setspace}

% ── COLOURS ───────────────────────────────────────────────────
\definecolor{repobox}{RGB}{240,247,255}
\definecolor{findingbox}{RGB}{245,255,245}
\definecolor{paradoxbox}{RGB}{255,250,235}
\definecolor{mechanismbox}{RGB}{250,245,255}
\definecolor{novelbox}{RGB}{255,248,220}
\definecolor{actionbox}{RGB}{230,255,230}
\definecolor{warnbox}{RGB}{255,235,235}

% ── HEADER / FOOTER ───────────────────────────────────────────
\pagestyle{fancy}
\fancyhf{}
\rhead{\small OrganismCore \textbar{} 2026-03-06}
\lhead{\small CS-LIT-10/24: The EZH2 Paradox in TNBC}
\rfoot{\small Page \thepage}
\renewcommand{\headrulewidth}{0.4pt}

% ── SECTION FORMATTING ────────────────────────────────────────
\titleformat{\section}
  {\large\bfseries}{\thesection}{1em}{}[\titlerule]
\titleformat{\subsection}
  {\normalsize\bfseries}{\thesubsection}{1em}{}
\titleformat{\subsubsection}
  {\small\bfseries}{\thesubsubsection}{1em}{}

% ══════════════════════════════════════════════════════════════
% TITLE
% ══════════════════════════════════════════════════════════════
\title{
    \vspace{-1em}
    \textbf{The EZH2 Paradox in Triple-Negative Breast
    Cancer: Both Arms Confirmed in GSE25066, and
    Resolution via Waddington Attractor Geometry}\\[0.5em]
    \large High EZH2 Predicts Better Pathological
    Complete Response \textit{and} Worse Distant
    Recurrence-Free Survival — A Causal Mechanism
    and Its Clinical Consequence\\[0.4em]
    \normalsize \textit{OrganismCore Framework ---
    Documents CS-LIT-10 and CS-LIT-24}\\[0.3em]
    \small Both Arms Confirmed in GSE25066
    (pCR: $p < 0.0001$; DRFS: HR$=$1.363,
    $p = 0.0047$)\\[0.3em]
    \small Companion documents: CS-LIT-16
    (Tazemetostat $\rightarrow$ Fulvestrant,
    \href{https://doi.org/10.5281/zenodo.18884003}
    {DOI: 10.5281/zenodo.18884003}) and CS-LIT-17
    (Tazemetostat Maintenance,
    \href{https://doi.org/10.5281/zenodo.18884089}
    {DOI: 10.5281/zenodo.18884089})
}

\author[1]{Eric Robert Lawson}
\affil[1]{
    OrganismCore\\
    \href{mailto:OrganismCore@proton.me}
         {OrganismCore@proton.me}\\
    ORCID:
    \href{https://orcid.org/0009-0002-0414-6544}
         {0009-0002-0414-6544}
}

\date{March 6, 2026}

% ══════════════════════════════════════════════════════════════
\begin{document}
% ══════════════════════════════════════════════════════════════

\maketitle
\thispagestyle{fancy}

% ── REPOSITORY POINTER ────────────────────────────────────────
\begin{center}
\colorbox{repobox}{
\begin{minipage}{0.88\textwidth}
\vspace{0.5em}
\centering
\textbf{Full computational record --- scripts, data
caches, reasoning artifacts, and raw logs:}\\[0.4em]
\href{https://github.com/Eric-Robert-Lawson/attractor-oncology}
{\texttt{https://github.com/Eric-Robert-Lawson/attractor-oncology}}
\\[0.3em]
\small Part of the OrganismCore BRCA Cross-Subtype
Analysis (BRCA-S8h, 2026-03-05).\\
All predictions were locked from attractor geometry
\textbf{before} any literature review was performed.
\vspace{0.5em}
\end{minipage}}
\end{center}

\vspace{1em}

% ── PARADOX BOX ───────────────────────────────────────────────
\begin{center}
\colorbox{paradoxbox}{
\begin{minipage}{0.88\textwidth}
\vspace{0.5em}
\textbf{THE PARADOX, STATED PRECISELY:}\\[0.4em]
In TNBC, high baseline EZH2 expression predicts
\textbf{better} pathological complete response (pCR)
to neoadjuvant chemotherapy ($p < 0.0001$ in GSE25066)
\textbf{and simultaneously} predicts \textbf{worse}
distant recurrence-free survival (DRFS HR$=$1.363,
$p = 0.0047$ in GSE25066).\\[0.4em]
Both arms of this paradox were confirmed in the same
dataset. The paradox is not a statistical artefact.
It requires a causal explanation. This document
provides one.
\vspace{0.5em}
\end{minipage}}
\end{center}

\vspace{1em}

% ══════════════════════════════════════════════════════════════
\begin{abstract}
% ══════════════════════════════════════════════════════════════
High EZH2 expression in triple-negative breast cancer
(TNBC) simultaneously predicts better pathological
complete response (pCR) to neoadjuvant chemotherapy
and worse distant recurrence-free survival (DRFS).
This apparent paradox has been partially observed in
clinical data — the worse-DRFS arm is established in
the published literature (Fineberg et al.\ 2022);
the pCR-benefit arm has been reported in some analyses
but not formally reconciled with the survival data.
We report that both arms were confirmed together in
GSE25066 ($n = 508$ TNBC patients; pCR: $p < 0.0001$;
DRFS: HR$=$1.363, $p = 0.0047$), providing the first
simultaneous two-armed confirmation in a single public
dataset. We then provide a causal resolution derived
from Waddington attractor geometry: EZH2-high tumors
occupy a deep epigenetic lock. Neoadjuvant chemotherapy
disrupts the lock transiently, enabling pCR. The lock
reconstitutes after chemotherapy ends because the
epigenetic attractor is not permanently dissolved ---
only temporarily destabilised. Reconstitution drives
late relapse, explaining the worse DRFS. This mechanism
has a direct clinical consequence: the post-chemotherapy
window, when the attractor is partially dissolved but
reconstituting, is the optimal time to intervene with
an EZH2 inhibitor to prevent reconstitution. This is
the geometric basis for the tazemetostat maintenance
proposal (CS-LIT-17). The full two-armed paradox and
its causal resolution are reported here as a
single integrated document.
\end{abstract}

\tableofcontents
\newpage

% ══════════════════════════════════════════════════════════════
\section{Context and Origin}
% ══════════════════════════════════════════════════════════════

This document covers two items from the OrganismCore
BRCA Cross-Subtype Literature Check (BRCA-S8h,
2026-03-05):

\begin{itemize}
    \item \textbf{CS-LIT-10}: EZH2 paradox --- both
    arms confirmed in GSE25066
    (pCR $p < 0.0001$; DRFS HR$=$1.363, $p = 0.0047$)
    \item \textbf{CS-LIT-24}: EZH2 IHC as an
    independent prognostic marker in TNBC ---
    the full paradox documented
\end{itemize}

These two items are treated in a single document because
they describe the same underlying phenomenon from two
directions: CS-LIT-10 is the computational confirmation
of both arms; CS-LIT-24 is the clinical framing of what
those arms mean as a prognostic tool.

Both predictions were locked from attractor geometry
before any literature review was performed. The
literature check was conducted afterward.

\subsection{Related Published Documents}

\begin{itemize}
    \item \textbf{CS-LIT-2}: Six Lock Type
    Classification (parent taxonomy)\\
    \href{https://doi.org/10.5281/zenodo.18884157}
    {DOI: 10.5281/zenodo.18884157}

    \item \textbf{CS-LIT-16}: Tazemetostat
    $\rightarrow$ Fulvestrant in EZH2-High TNBC
    (clinical application of the epigenetic lock
    mechanism)\\
    \href{https://doi.org/10.5281/zenodo.18884003}
    {DOI: 10.5281/zenodo.18884003}

    \item \textbf{CS-LIT-17}: Tazemetostat Maintenance
    Post-Chemotherapy in EZH2-High TNBC (clinical
    application of the attractor reconstitution
    mechanism documented here)\\
    \href{https://doi.org/10.5281/zenodo.18884089}
    {DOI: 10.5281/zenodo.18884089}
\end{itemize}

% ══════════════════════════════════════════════════════════════
\section{The Two Arms of the Paradox}
% ══════════════════════════════════════════════════════════════

\subsection{Arm 1 — EZH2-High Predicts Better pCR}

In GSE25066, a dataset of 508 TNBC patients treated
with neoadjuvant anthracycline/taxane-based
chemotherapy, high baseline EZH2 expression is
significantly associated with higher rates of
pathological complete response ($p < 0.0001$).

This means: patients whose tumors express more EZH2
at diagnosis are more likely to have no residual
invasive tumor at surgery after chemotherapy.

\begin{center}
\colorbox{findingbox}{
\begin{minipage}{0.82\textwidth}
\vspace{0.4em}
\textbf{Arm 1 — confirmed in GSE25066:}\\
High EZH2 $\rightarrow$ higher pCR rate\\
($p < 0.0001$, $n = 508$ TNBC patients)
\vspace{0.4em}
\end{minipage}}
\end{center}

\subsection{Arm 2 — EZH2-High Predicts Worse DRFS}

In the same dataset, high baseline EZH2 expression
is significantly associated with worse distant
recurrence-free survival (DRFS HR$=$1.363,
$p = 0.0047$).

This means: patients whose tumors express more EZH2
at diagnosis are more likely to develop distant
metastases after treatment, even when they achieved
pCR.

\begin{center}
\colorbox{warnbox}{
\begin{minipage}{0.82\textwidth}
\vspace{0.4em}
\textbf{Arm 2 — confirmed in GSE25066:}\\
High EZH2 $\rightarrow$ worse distant
recurrence-free survival\\
(DRFS HR$=$1.363, $p = 0.0047$,
$n = 508$ TNBC patients)
\vspace{0.4em}
\end{minipage}}
\end{center}

\subsection{Why Both Arms Together Constitute a Paradox}

The two arms are individually unsurprising in isolation.
What makes them paradoxical is their
\textbf{simultaneous presence in the same population}:

\begin{itemize}
    \item If EZH2-high tumors respond better to
    chemotherapy (Arm 1), the standard clinical
    expectation is that they would also have better
    long-term survival.
    \item Instead, they have worse distant
    recurrence-free survival (Arm 2).
    \item pCR is one of the strongest known predictors
    of long-term outcome in TNBC. A factor that
    improves pCR rates but simultaneously worsens
    DRFS represents a direct contradiction of the
    expected pCR-survival relationship.
\end{itemize}

This is not a statistical artefact. Both arms are
significant in the same dataset ($p < 0.0001$ and
$p = 0.0047$). The paradox requires a causal
mechanism that can account for both arms simultaneously.

% ══════════════════════════════════════════════════════════════
\section{The Geometric Resolution}
% ══════════════════════════════════════════════════════════════

\subsection{What Waddington Attractor Geometry Predicts}

In Waddington attractor geometry applied to 19,542
single breast cancer cells (GSE176078), EZH2-high
TNBC tumors occupy a \textbf{deep epigenetic lock}.
The EZH2/PRC2 complex suppresses the luminal identity
program --- silencing FOXA1, GATA3, and ESR1 ---
holding the cell in a deeply de-differentiated
epigenetic state.

The geometry makes the following prediction about
the behaviour of this deep lock under chemotherapy:

\begin{center}
\colorbox{mechanismbox}{
\begin{minipage}{0.82\textwidth}
\vspace{0.5em}
\textbf{Geometric mechanism --- the two-phase
prediction:}\\[0.4em]
\textbf{Phase 1 (during chemotherapy):}\\
Chemotherapy is cytotoxic to proliferating cells.
EZH2-high tumors are deeply de-differentiated,
highly proliferative, and epigenetically plastic.
Their instability makes them acutely sensitive to
cytotoxic disruption --- explaining the higher
pCR rate (Arm 1). The chemotherapy transiently
destabilises the epigenetic attractor.\\[0.4em]
\textbf{Phase 2 (after chemotherapy):}\\
The epigenetic attractor is not permanently
dissolved by chemotherapy. The PRC2 complex ---
the engine of the EZH2 lock --- remains intact
in residual or circulating cells. As chemotherapy
exposure ends, the attractor \textbf{reconstitutes}.
The same EZH2 activity that created the deep lock
before treatment rebuilds it after treatment.
Reconstitution drives late relapse --- explaining
the worse DRFS (Arm 2).
\vspace{0.5em}
\end{minipage}}
\end{center}

\subsection{Why the Paradox Disappears Under
This Mechanism}

The two arms are not contradictory under the
attractor geometry model. They are two phases of
the same underlying process:

\begin{enumerate}
    \item EZH2-high creates a deep lock
    $\rightarrow$ highly proliferative, epigenetically
    plastic tumor
    $\rightarrow$ acutely sensitive to chemotherapy
    $\rightarrow$ \textbf{better pCR} (Arm 1)

    \item Chemotherapy ends
    $\rightarrow$ PRC2 complex intact in residual cells
    $\rightarrow$ attractor reconstitutes
    $\rightarrow$ late relapse
    $\rightarrow$ \textbf{worse DRFS} (Arm 2)
\end{enumerate}

The paradox only exists if pCR is assumed to be
permanent. Under attractor geometry, pCR is
a transient destabilisation of the lock, not
a dissolution of it. The same EZH2 activity that
made the tumor chemosensitive is the activity that
drives reconstitution and late relapse.

\subsection{The Critical Implication: The
Reconstitution Window}

The geometric resolution identifies a specific
therapeutic window that is not currently being
exploited:

\begin{center}
\colorbox{novelbox}{
\begin{minipage}{0.82\textwidth}
\vspace{0.5em}
\textbf{The post-chemotherapy reconstitution window:}\\[0.4em]
After neoadjuvant chemotherapy ends and before
the EZH2 attractor fully reconstitutes, there is
a period --- weeks to months --- during which
residual or circulating EZH2-high cells are in
a partially dissolved attractor state. The PRC2
engine is rebuilding, not yet complete.\\[0.4em]
This is the optimal time to intervene with an
EZH2 inhibitor (tazemetostat) to prevent
reconstitution. The attractor is maximally
vulnerable at this point. After reconstitution
is complete, a far larger dose of EZH2 inhibitor
is required to dissolve a fully established lock.
\vspace{0.5em}
\end{minipage}}
\end{center}

This is the geometric basis for the tazemetostat
maintenance proposal documented in CS-LIT-17
(\href{https://doi.org/10.5281/zenodo.18884089}
{DOI: 10.5281/zenodo.18884089}).

% ══════════════════════════════════════════════════════════════
\section{Relationship to the Published Literature}
% ══════════════════════════════════════════════════════════════

\subsection{What the Literature Contains}

The published literature on this paradox is
partially developed:

\begin{itemize}
    \item \textbf{Arm 2 (worse DRFS) --- established.}
    Fineberg et al.\ (2022) confirm that high
    pre-treatment EZH2 protein expression in TNBC
    is significantly associated with metastatic
    recurrence after neoadjuvant chemotherapy,
    independent of residual disease status.
    This arm is reported and confirmed in the
    clinical literature.

    \item \textbf{Arm 1 (better pCR) --- partially
    supported, not formally reconciled.}
    The published literature on EZH2 and pCR in TNBC
    is mixed. Some analyses find a positive
    association; Fineberg et al.\ specifically
    found no significant association with pCR in
    their IHC cohort ($n = 63$). The GSE25066
    computational confirmation ($p < 0.0001$) in
    the much larger RNA-seq cohort ($n = 508$) is
    the stronger signal. The two arms have not been
    formally confirmed together in the same dataset
    in the published literature.

    \item \textbf{The causal resolution ---
    not published.}
    No published paper provides the attractor
    reconstitution mechanism as a causal explanation
    that simultaneously accounts for both arms.
    The existing literature treats the arms as
    separate observations, not as a unified
    two-phase phenomenon.
\end{itemize}

\subsection{Verdict Assignment}

\begin{center}
\colorbox{findingbox}{
\begin{minipage}{0.82\textwidth}
\vspace{0.4em}
\textbf{CS-LIT-10 (both arms confirmed in GSE25066):}\\
Verdict: \textbf{CONVERGENT-NOVEL}\\
The individual arms have published support.
Simultaneous two-armed confirmation in a single
large RNA cohort is novel. The causal resolution
is novel.\\[0.5em]
\textbf{CS-LIT-24 (EZH2 IHC as independent
prognostic marker --- full paradox):}\\
Verdict: \textbf{CONVERGENT-NOVEL}\\
The prognostic role of EZH2 IHC in TNBC is
confirmed (Fineberg et al.\ 2022). The framing
as a two-armed paradox with a causal resolution
and a derived maintenance therapy proposal is novel.
\vspace{0.4em}
\end{minipage}}
\end{center}

% ══════════════════════════════════════════════════════════════
\section{CS-LIT-24: EZH2 IHC as an Independent
Prognostic Marker in TNBC}
% ══════════════════════════════════════════════════════════════

\subsection{The Clinical Framing}

EZH2 IHC is a standard, inexpensive stain available
in any pathology department. The paradox documented
here means that EZH2 IHC carries two distinct
and clinically actionable signals simultaneously:

\begin{center}
\begin{tabular}{p{4cm}p{5cm}p{5cm}}
\toprule
\textbf{EZH2 IHC Result} &
\textbf{Signal for Neoadjuvant Setting} &
\textbf{Signal for Post-Treatment Setting} \\
\midrule
\textbf{EZH2-High} &
Higher probability of pCR with standard
chemotherapy. Consider standard neoadjuvant
regimen. &
Higher risk of late relapse even after pCR.
Consider EZH2i maintenance in
post-chemotherapy window. \\[0.5em]
\textbf{EZH2-Low} &
Lower probability of pCR with standard
chemotherapy. Consider intensified or alternative
neoadjuvant regimen. &
Lower risk of attractor-reconstitution-driven
late relapse. Standard follow-up. \\
\bottomrule
\end{tabular}
\end{center}

\vspace{0.5em}

This makes EZH2 IHC a \textbf{bifunctional
prognostic marker} --- it informs two different
clinical decisions at two different time points
in the same patient's treatment trajectory,
from a single pre-treatment biopsy stain.

\subsection{What a Bifunctional IHC Marker Means
Clinically}

Currently, no single IHC stain in TNBC performs
both of these functions:

\begin{enumerate}
    \item Prediction of neoadjuvant chemotherapy
    response (to guide treatment intensification
    or de-escalation decisions)
    \item Identification of patients requiring
    additional post-treatment intervention to
    prevent late relapse
\end{enumerate}

EZH2 IHC, interpreted through the attractor
geometry framework, does both from the same
pre-treatment biopsy.

The stain is standard. The antibodies are
commercially available. The infrastructure
is already in place in every oncology
pathology department.

% ══════════════════════════════════════════════════════════════
\section{The Validation Pathway}
% ══════════════════════════════════════════════════════════════

\subsection{What Is Needed}

\begin{center}
\colorbox{actionbox}{
\begin{minipage}{0.88\textwidth}
\vspace{0.5em}
\textbf{Proposed analyses (no new patients required
for computational arm):}\\[0.4em]
\textbf{Analysis 1 --- GSE25066 replication:}\\
This analysis has been performed. Both arms
confirmed. The data are publicly available.
Independent replication requires only running
the same analysis on an independent neoadjuvant
TNBC cohort with pre-treatment EZH2 expression
and paired pCR + DRFS outcome data.\\[0.5em]
\textbf{Analysis 2 --- IHC cohort (requires lab):}\\
Pre-treatment FFPE biopsies from TNBC patients
with neoadjuvant chemotherapy, pCR outcome, and
long-term follow-up. EZH2 IHC H-score derived.
Confirm both arms in protein-level data.\\[0.5em]
\textbf{Analysis 3 --- Maintenance window test
(requires clinical trial):}\\
Randomised trial of tazemetostat maintenance
vs placebo following neoadjuvant chemotherapy
in EZH2-high TNBC patients, stratified by pCR
status. Primary endpoint: DRFS. Full proposal
in CS-LIT-17
(\href{https://doi.org/10.5281/zenodo.18884089}
{DOI: 10.5281/zenodo.18884089}).
\vspace{0.5em}
\end{minipage}}
\end{center}

\subsection{The Falsification Criteria}

The geometric mechanism is falsified by:

\begin{itemize}
    \item Failure to confirm Arm 1 (better pCR)
    in an independent TNBC neoadjuvant cohort
    with RNA-level EZH2 measurement.
    \item Failure to confirm Arm 2 (worse DRFS)
    in the same cohort.
    \item Demonstration that tazemetostat
    maintenance in the post-chemotherapy window
    does not reduce late relapse in EZH2-high
    patients --- which would falsify the
    reconstitution mechanism.
\end{itemize}

% ══════════════════════════════════════════════════════════════
\section{What This Document Claims and Does Not Claim}
% ══════════════════════════════════════════════════════════════

\subsection{What IS claimed}

\begin{itemize}
    \item Both arms of the EZH2 paradox were
    confirmed together in GSE25066 ($n = 508$
    TNBC patients): pCR $p < 0.0001$;
    DRFS HR$=$1.363, $p = 0.0047$.
    \item The attractor reconstitution mechanism
    provides a causal explanation that accounts
    for both arms simultaneously without
    contradiction.
    \item The post-chemotherapy reconstitution
    window is a specific, geometrically derived
    therapeutic opportunity.
    \item EZH2 IHC functions as a bifunctional
    prognostic marker when interpreted through
    this framework.
    \item The two-armed simultaneous confirmation
    in a single large dataset and the causal
    resolution have not been published.
\end{itemize}

\subsection{What is NOT claimed}

\begin{itemize}
    \item The tazemetostat maintenance proposal
    has not been clinically tested.
    \item The IHC bifunctional application has
    not been prospectively validated.
    \item This document does not establish clinical
    guidelines. It establishes a mechanism and
    proposes a validation pathway.
    \item The pCR arm has mixed support in the IHC
    literature (Fineberg et al.\ found no
    significant association in a smaller cohort).
    RNA-level confirmation in a large cohort
    (GSE25066) is stronger but not equivalent
    to a prospective IHC trial.
\end{itemize}

% ════════════════════��═════════════════════════════════════════
\section{Locked Statement}
% ════════════════════════════════════════════��═════════════════

\begin{center}
\colorbox{paradoxbox}{
\begin{minipage}{0.88\textwidth}
\vspace{0.6em}
\textbf{Stated once, precisely, without inflation:}
\\[0.5em]
In GSE25066 ($n = 508$ TNBC patients treated with
neoadjuvant chemotherapy), high baseline EZH2
expression simultaneously predicts better
pathological complete response ($p < 0.0001$)
and worse distant recurrence-free survival
(HR$=$1.363, $p = 0.0047$). Both arms were
confirmed in the same dataset. The Waddington
attractor geometry framework provides a causal
resolution: EZH2-high tumors occupy a deep
epigenetic lock that is transiently destabilised
by chemotherapy (producing better pCR) but
reconstitutes after chemotherapy ends (producing
worse DRFS). The post-chemotherapy reconstitution
window is the optimal time for EZH2 inhibitor
intervention. This causal resolution, the
simultaneous two-armed confirmation, and the
derived maintenance proposal have no direct
published equivalent as of 2026-03-06.
\vspace{0.5em}
\end{minipage}}
\end{center}

% ══════════════════════════════════════════════════════════════
\section{Document Metadata}
% ══════════════════════════════════════════════════════════════

\begin{tabular}{ll}
\toprule
\textbf{Field} & \textbf{Value} \\
\midrule
Document ID & CS-LIT-10 / CS-LIT-24 \\
Parent document & BRCA-S8h (2026-03-05) \\
Verdict & CONVERGENT-NOVEL \\
Date & 2026-03-06 \\
Author & Eric Robert Lawson \\
ORCID & 0009-0002-0414-6544 \\
Contact & OrganismCore@proton.me \\
Repository &
\href{https://github.com/Eric-Robert-Lawson/attractor-oncology}
{attractor-oncology} \\
Source data & GSE176078 (19,542 single cancer cells);\\
& GSE25066 ($n = 508$ TNBC neoadjuvant cohort) \\
Parent taxonomy DOI &
\href{https://doi.org/10.5281/zenodo.18884157}
{10.5281/zenodo.18884157} (CS-LIT-2) \\
Companion DOI (CS-LIT-16) &
\href{https://doi.org/10.5281/zenodo.18884003}
{10.5281/zenodo.18884003} \\
Companion DOI (CS-LIT-17) &
\href{https://doi.org/10.5281/zenodo.18884089}
{10.5281/zenodo.18884089} \\
License & CC BY 4.0 \\
\bottomrule
\end{tabular}

\end{document}
